\documentclass[10pt,a4paper]{amsart}
\usepackage[margin=19mm]{geometry}
\usepackage{amsmath,amssymb,mathtools}
\usepackage[scr=rsfs,cal=euler]{mathalfa}
\usepackage{xfrac}
\usepackage[unicode,hidelinks,bookmarks=false]{hyperref}
\newcommand{\ie}{i.e.\ }
\newcommand{\cf}{cf.\ }
\newcommand{\resp}{respectively\ }
\newcommand{\Subsection}{\subsection}
\providecommand{\nobreakdash}{\nobreak\hbox{-}}
\setlength{\emergencystretch}{2em}
\multlinegap0pt
\usepackage{xcolor}
\newtheorem{lemma}{Lemma}
\usepackage{cleveref}
\crefname{lemma}{Lemma}{Lemmas}
\newcommand{\F}{\mathbb{F}}
\def\U{\mathcal{U}}
\def\V{\mathcal{V}} 
\newcommand{\Z}{\mathbb{Z}}
\title{Improved Subexponential Upper Bounds for $3$-Restricted Matching Vector Families}
\author{Sidhant Saraogi}
\date{Source: arXiv:2608.27859v1 (CC BY 4.0)}
\begin{document}
\maketitle

\noindent\emph{Source and licence.} Improved Subexponential Upper Bounds for $3$-Restricted Matching Vector Families. Version arXiv:2608.27859v1 (CC BY 4.0), distributed by arXiv under the Creative Commons Attribution 4.0 International licence, http://creativecommons.org/licenses/by/4.0/, as recorded in the arXiv licence metadata for this exact version and shown on the abstract page https://arxiv.org/abs/2608.27859v1. Full licence notice: \texttt{licenses/arxiv-2608.27859-CC-BY-4.0.txt}.\par\smallskip

\noindent\textit{Editorial scope.} Counted unit: the article's lemma lem:equal-sum-moonflower (source0.tex lines 401--410). The article's complete original proof of it is delivered with it (source0.tex lines 412--528, one \texttt{proof} environment of 117 lines). No mathematics is written, repaired or re-derived: the statement and the proof are the article's own lines, in the article's own notation, and no retained line is substituted or re-flowed.  Retained uncounted as the source-local context the proof needs: lines 294--309; lines 340--355.  Nothing was omitted from any retained span, so the package carries no omission record.  No retained line contains a citation, so the wrapper carries no bibliography and no bibliography entry.  No retained line contains a figure environment, an image-inclusion command or drawing code, so no image is delivered and the package declares no asset dependency.  Editorial formatting only, in addition: an AMS \texttt{amsart} preamble that restates the article's own declarations and macros used by the retained lines, namely the environments \texttt{lemma} on separate counters, together with the standard packages those lines need.  The retained environments print in this document as Lemma 1 in order of appearance, whereas the article prints the same items with the numbering of its own full document.  The complete native source of the article as distributed in the arXiv e-print for this exact version is bundled unchanged as \texttt{sources/208744/source0.tex}.
\begin{lemma} \label{lem:coeff_rank}
Let $R=\mathbb Z_{q^b}$ where $q$ is a prime and $b \geq 1$. Let
\[
X\in R^{M\times n},
\qquad
Y\in R^{n\times N}.
\]
Suppose every entry of $XY$ lies in $q^{b-1}R$. Define $L\in\mathbb F_q^{M\times N}$ by
\[
L_{ij}=\tau_q((XY)_{ij}).
\]
Then
\[
\operatorname{rank}_{\mathbb F_q}(L)\le n.
\]
\end{lemma}

\begin{lemma} \label{lem:polynomial_rank}
Let $B_1,\ldots,B_d$ be $M\times N$ matrices over a field $\mathbb F$, and let
\[
\operatorname{rank}_{\mathbb F}(B_h)\le n
\qquad
\text{for every }h\in[d].
\]
Let $f\in\mathbb F[z_1,\ldots,z_d]$ be multilinear, and define $P$ entry-wise by
\[
P_{ij}=f((B_1)_{ij},\ldots,(B_d)_{ij}).
\]
Then
\[
\operatorname{rank}_{\mathbb F}(P)\le (n+1)^d.
\]
\end{lemma}

\begin{lemma}[Equal-sum moonflower bound]\label{lem:equal-sum-moonflower}
Let $\U,\V$ be a typical $3$-restricted MVF modulo $m=p^aq^b$. For every $w\in\Z_m^n$ and every $d\ge 1$, if
\[
A_1,\ldots,A_L \in C_w^{(d)}
\]
is a moonflower, then
\[
L\le d\left((n+1)^d+1\right).
\]
\end{lemma}

\begin{proof}
Fix $w\in\Z_m^n$, $d\ge 1$, and a moonflower
\[
A_1,\ldots,A_L \in C_w^{(d)}.
\]
The proof has three main components. We begin with the full matrix of inner products
\[
M_{ij}:=\langle u_j,v_i\rangle\in\Z_m .
\]
The diagonal entries of $M$ are $0$, while every off-diagonal entry lies in
$S = \{\alpha,\beta\}$.

\paragraph{Moving from $\Z_m$ to $\F_q$.} First, we reduce all inner products modulo $q^b$. Since
\[
\alpha\in q^{b-1}\mathbb Z_{q^b},
\qquad
\beta\equiv 0\pmod {q^b},
\]
and the diagonal entries are $0$, all entries lie in $q^{b-1}\mathbb Z_{q^b}$. Define
\[
Q_{ij}
=
\tau_q(\langle u_j,v_i\rangle).
\]
By \cref{lem:coeff_rank}, we know that
\[
\operatorname{rank}_{\mathbb F_q}(Q)\le n.
\]
Moreover, since $\alpha \not\equiv 0\pmod{q^b}$, there is a nonzero $\gamma\in\mathbb F_q$ such that
\[
Q_{ij}
=
\gamma\cdot \mathbf 1[\langle u_j,v_i\rangle=\alpha].
\]
Let $\Pi = \gamma^{-1} Q$, then we have $\operatorname{rank}_{\mathbb F_q}(\Pi) = \operatorname{rank}_{\mathbb F_q}(Q) \le n$. Note that $\Pi_{ij} = 1[\langle u_j,v_i\rangle=\alpha].$ $\Pi$ is a low-rank indicator matrix for the residue $\alpha$. 

\paragraph{From elements to sums.} Now, for each $h\in[L]$, choose a unique element
\[
i_h\in A_h,
\qquad
i_h\notin A_t\quad\text{for all }t\ne h. 
\]
We will use
$\Pi$ to build a new matrix indexed by the sets $A_1,\ldots,A_L$, capturing the
interaction between the sums $u_{A_t}$ and the unique element $v_{i_h}$.
Define
\[
\ell_h
=
\left|\{j\in A_h:\langle u_j,v_{i_h}\rangle=\alpha\}\right|.
\]
Since $i_h\in A_h$ gives the diagonal inner product $0$, we have
$\ell_h\in\{0,\ldots,d-1\}$. By the pigeonhole principle, some value $\ell$ occurs for at least $L/d$ of the sets. Restricting to those sets and relabeling, we obtain a sub-moonflower $A_1,\ldots,A_{L'}$ with $L'\ge L/d$ and $\ell_h=\ell$ throughout.

For $h\ne t$, let
\[
r_{ht}
=
\left|\{j\in A_t:\langle u_j,v_{i_h}\rangle=\alpha\}\right|.
\]
Since $i_h\notin A_t$, all terms in
$\langle u_{A_t},v_{i_h}\rangle$ are off-diagonal. Since
$u_{A_h}=u_{A_t}$, taking their respective inner products with $v_{i_h}$ gives
\[
r_{ht}\alpha+(d-r_{ht})\beta
\equiv
\ell\alpha+(d-1-\ell)\beta
\pmod m.
\]
Thus $r_{ht}$ must lie in the set
\[
T_\ell
=
\left\{
r\in\{0,\ldots,d\}:
r\alpha+(d-r)\beta
\equiv
\ell\alpha+(d-1-\ell)\beta
\pmod m
\right\}.
\]
Crucially, $\ell\notin T_\ell,$
since $r=\ell$ would force $\beta\equiv0\pmod m$. Therefore, along the diagonal block pair $(h,h)$ the row indexed by $i_h$
sees exactly $\ell$ occurrences of the residue $\alpha$, while along every
off-diagonal block pair $(h,t)$ it sees a number of occurrences belonging to
$T_\ell$, a set that avoids $\ell$.

\paragraph{Designing the polynomial method matrix.}  We can encode this with a function $g:\{0,1\}^d\to\mathbb F_q$ such that
\[
g(z)=0\quad\text{if }|z|=\ell,
\qquad
g(z)=1\quad\text{if }|z|\in T_\ell.
\] 
Let $f\in\mathbb F_q[z_1,\ldots,z_d]$ be the unique multilinear polynomial agreeing with $g$ on the Boolean cube. Fix orderings $A_t=\{j_{t,1},\ldots,j_{t,d}\}$ and define
\[
P_{ht}:=
f(\Pi_{i_h,j_{t,1}},\ldots,\Pi_{i_h,j_{t,d}}).
\]
By construction, the input has Hamming weight $\ell$ when $h=t$, and
Hamming weight in $T_\ell$ when $h\ne t$. Hence $P=J-I$, so
\begin{equation} \label{eqn:lower}
L/d - 1\leq L'-1\le \operatorname{rank}_{\mathbb F_q}(P).
\end{equation}
Now, for $r\in[d]$, define
\[
(B_r)_{ht}:=\Pi_{i_h,j_{t,r}}.
\]
Each $B_r$ is obtained from $\Pi$ by selecting rows and columns, possibly with
repetitions, so \\ $\operatorname{rank}_{\mathbb F_q}(B_r)\le n$. Since $P$ is
obtained by applying $f$ entrywise to $B_1,\ldots,B_d$, \cref{lem:polynomial_rank}
gives
\begin{equation} \label{eqn:upper}
\operatorname{rank}_{\mathbb F_q}(P)\le (n+1)^d.
\end{equation}

Combining \cref{eqn:lower} and \cref{eqn:upper} completes our proof.
\end{proof}

\end{document}
